\begin{titlepage}
    \centering
    \vspace*{0.2cm}


    \includesvg[width=0.17\textwidth]{isi-logo.svg}\par
    \vspace{0.4cm}
    {\Large \textsc{Indian Statistical Institute, Kolkata}}\par
    \vspace{0.1cm}
    {\large \textsc{Bachelor of Statistics}}\par

    \vspace{0.5cm}
    \rule{\textwidth}{0.8pt}\par
    \vspace{0.4cm}


    {\large \itshape Economic Statistics and Demography \textbullet{} Project Report}\par
    \vspace{0.5cm}
    {\LARGE \bfseries \textsc{Storytelling with Data:}}\par
    \vspace{0.25cm}
    {\Large \textsc{A Comparative Demographic and Economic Study of\\ Sikkim and Rajasthan}}\par

    \vspace{0.4cm}
    \rule{\textwidth}{0.8pt}\par
    \vspace{0.6cm}
    {\large
    \begin{tabular}{c c}
        Shounak Kar (BS2439)    & Sayan Kar (BS2436) \\[0.2cm]
        Ahan Chakraborty (BS2402) & Arnab Sanyal (BS2411)
    \end{tabular}\par}

    \vspace{0.5cm}
    {\large \textbf{Advisor:} Prof. Kajori Banerjee}\par
    \vspace{0.2cm}
    {\large 12\textsuperscript{th} October 2026}\par

    \vspace{0.8cm}
    
    \begin{minipage}{0.88\textwidth}
        \centering
        {\large \textbf{Abstract}}\par
        \vspace{1cm}
        \justifying
        \small
        \noindent
        

This study examines demographic change, gender differences in health, household consumption and health-financing coverage in Rajasthan and Sikkim using survey-weighted microdata from the NSS 80th Round Health Survey (2025). The two states differ substantially in age structure, household size, marriage patterns and women's educational attainment. Age-standardised chronic-ailment prevalence is broadly similar across states and between men and women. However, hospitalisation patterns reveal an important gender contrast: in Rajasthan, women aged outside 15–49 have lower hospitalisation rates than men, despite comparable reported morbidity, a finding that warrants further investigation. Median monthly per-capita consumption expenditure is higher in Sikkim, although overall consumption inequality is similar in the two states. The largest contrast concerns reported health-financing coverage, estimated at 77.2\% in Rajasthan and 18.6\% in Sikkim. Coverage in Rajasthan is distributed relatively evenly across consumption groups, whereas Sikkim exhibits a more pronounced, though statistically imprecise, concentration among richer households. These findings illustrate important differences in demographic structure, gender-related outcomes and reported health-scheme coverage, while highlighting the limitations of cross-sectional data for drawing causal conclusions about healthcare access and policy effectiveness.\par
    \end{minipage}

    \vfill
\end{titlepage}